\section*{अध्याय \ref{ch:b1:crystals-ionic-covalent} --- क्रिस्टल II: आयनिक, सहसंयोजी और आण्विक ठोस}

\begin{solution}{exo:b1:crystals-ionic-covalent:1}
ऋणायन: $8 \times \frac18 + 6 \times \frac12 = 4$; धनायन: $12 \times
\frac14 + 1 = 4$। हर आयन के दूसरे प्रकार के 6 पड़ोसी हैं (6:6)। चार
धनायन और चार ऋणायन: \ce{NaCl}।
\end{solution}

\begin{solution}{exo:b1:crystals-ionic-covalent:2}
मुख्य विकर्ण का आधा: $a\sqrt3/2 = 412.3 \times 0.866 = \qty{357.1}{pm}$।
एक \ce{Cs+} और $8 \times \frac18 = 1$ \ce{Cl-}: एक सूत्र इकाई।
\end{solution}

\begin{solution}{exo:b1:crystals-ionic-covalent:3}
$x = 102/181 = 0.56$, 0.414 और 0.732 के बीच: अष्टफलकीय उपसहसंयोजन
(6), \omterm{def:b1:crystals-ionic-covalent:structure-type}{सैंधव-लवण प्रकार}।
\end{solution}

\begin{solution}{exo:b1:crystals-ionic-covalent:4}
धात्विक: कॉपर। आयनिक: सोडियम क्लोराइड, कैल्सियम फ़्लुओराइड। सहसंयोजी:
हीरा, क्वार्ट्ज़। आण्विक: डाइआयोडीन, बर्फ़।
\end{solution}

\begin{solution}{exo:b1:crystals-ionic-covalent:5}
कोर के अनुदिश ऋणायन--धनायन--ऋणायन संपर्क में: $a = 2(r_+ + r_-)$।
फलक-विकर्ण पर दो ऋणायन दूरी $a\sqrt2/2$ पर हैं; वे एक-दूसरे पर
चढ़ने नहीं चाहिए: $a\sqrt2/2 \ge 2r_-$, इसलिए $(r_+ + r_-)\sqrt2 \ge 2r_-$, अर्थात्
$r_+/r_- \ge \sqrt2 - 1$।
\end{solution}

\begin{solution}{exo:b1:crystals-ionic-covalent:6}
$\rho = 4 \times 58.44/(6.022 \times 10^{23} \times (5.641 \times
10^{-8})^3) = \qty{2.16}{g/cm^3}$, जबकि मापा गया \qty{2.17}{g/cm^3}।
\end{solution}

\begin{solution}{exo:b1:crystals-ionic-covalent:7}
$\rho = 168.36/(6.022 \times 10^{23} \times (4.123 \times 10^{-8})^3) =
\qty{3.99}{g/cm^3}$, जो मापा गया मान ही है।
\end{solution}

\begin{solution}{exo:b1:crystals-ionic-covalent:8}
\ce{Zn-S} $= a\sqrt3/4 = \qty{234.2}{pm}$;
$\rho = 4 \times 97.44/(6.022 \times 10^{23} \times (5.409 \times
10^{-8})^3) = \qty{4.09}{g/cm^3}$, प्राकृतिक स्फेलेराइट के मापे गए
\qty{4.04}{g/cm^3} के पास।
\end{solution}

\begin{solution}{exo:b1:crystals-ionic-covalent:9}
\ce{C-C} $= 245.6/\sqrt3 = \qty{141.8}{pm}$; परतों के बीच $669.6/2 =
\qty{334.8}{pm}$। कोष्ठिका का आयतन $\frac{\sqrt3}{2}a^2c = \qty{3.498e-23}{cm^3}$,
$\rho = 4 \times 12.01/(6.022 \times 10^{23} \times 3.498 \times 10^{-23})
= \qty{2.28}{g/cm^3}$, जबकि हीरे के लिए \qty{3.52}{g/cm^3}: परतें
एक-दूसरे से दूर हैं और केवल लंदन बल उन्हें जोड़े रखते हैं।
\end{solution}

\begin{solution}{exo:b1:crystals-ionic-covalent:10}
\cref{prop:b1:crystals-ionic-covalent:diamond} की उपपत्ति देखिए:
$C = \pi\sqrt3/16 = 0.34$। कठोरता \omterm{def:b1:lewis-resonance-vsepr:lewis}{सहसंयोजी आबंधों} की प्रबलता और
त्रिविमीय दृढ़ता से आती है, संकुलन से नहीं: हर
परमाणु के केवल चार पड़ोसी हैं, निश्चित कोणों पर, जिससे बहुत ख़ाली
जगह बचती है।
\end{solution}

\begin{solution}{exo:b1:crystals-ionic-covalent:11}
\ce{Si-Si} $= a\sqrt3/4 = \qty{235.2}{pm}$, \omterm{def:b1:periodicity:radius}{सहसंयोजी त्रिज्या}
\qty{117.6}{pm}; $\rho = 8 \times 28.09/(6.022 \times 10^{23} \times
(5.431 \times 10^{-8})^3) = \qty{2.33}{g/cm^3}$, जो मापा गया मान ही है।
\end{solution}

\begin{solution}{exo:b1:crystals-ionic-covalent:12}
$x = 152/181 = 0.84 > 0.732$: नियम \omterm{def:b1:crystals-ionic-covalent:structure-type}{सीज़ियम क्लोराइड प्रकार} का पूर्वानुमान देता है।
सैंधव-लवण संरचना में \ce{Rb-Cl} $= a/2 = \qty{329.1}{pm}$, जो
$152 + 181 = \qty{333}{pm}$ के पास है: प्रेक्षित संरचना संगत है।
नियम आयनों को कठोर गोले मानता है और उनकी \omterm{def:b1:periodicity:polarisability}{ध्रुवणीयता} तथा दोनों संरचनाओं के बीच
ऊर्जा के छोटे अंतर को अनदेखा करता है (रुबिडियम
क्लोराइड दाब में \omterm{def:b1:crystals-ionic-covalent:structure-type}{सीज़ियम क्लोराइड प्रकार} अपना लेता है); यह एक मार्गदर्शक है,
कोई अटल नियम नहीं।
\end{solution}

\begin{solution}{pb:b1:crystals-ionic-covalent:1}
\textbf{1.} 8 कोनों और 6 फलक-केंद्रों पर: $8 \times \frac18 + 6
\times \frac12 = 4$ आयन।
\textbf{2.} 8 \omterm{def:b1:crystals-metals:site}{चतुष्फलकीय स्थल}, सभी भरे हुए: 8 \ce{F-}, 4
\ce{Ca^{2+}} के लिए, अनुपात 1:2, \ce{CaF2}।
\textbf{3.} \ce{F-}: चतुष्फलक के कोनों पर 4 \ce{Ca^{2+}}।
\ce{Ca^{2+}}: घन के कोनों पर 8 \ce{F-}।
\textbf{4.} $a\sqrt3/4 = 546.3 \times 0.4330 = \qty{236.6}{pm}$।
\textbf{5.} $112 + 131 = \qty{243}{pm}$: 3\,\% के भीतर।
\textbf{6.} \ce{F-} आयन कोर $a/2 =
\qty{273.2}{pm}$ वाली सरल घनीय व्यवस्था बनाते हैं, जो $2 \times 131 = \qty{262}{pm}$ से थोड़ा अधिक है: वे
लगभग छूते हैं।
\textbf{7.} $x = 112/131 = 0.85$।
\textbf{8.} $x > 0.732$: उपसहसंयोजन 8।
\textbf{9.} \omterm{def:b1:crystals-ionic-covalent:structure-type}{सीज़ियम क्लोराइड प्रकार} में धनायन और ऋणायन बराबर संख्या में होते हैं;
दोगुने ऋणायनों के साथ ऋणायनों के घनों में से केवल आधे में धनायन आ सकता है।
\textbf{10.} कोष्ठिका के 8 \ce{F-} कोर
$a/2$ वाले 8 छोटे घनों का एक घन बनाते हैं; \ce{Ca^{2+}} आयन उनमें से 4 के केंद्रों पर,
एकांतर क्रम में, बैठते हैं, जैसे त्रिविमीय शतरंज की बिसात के काले खाने।
\textbf{11.} $C = \frac43\pi(4 \times 112^3 + 8 \times 131^3)/546.3^3 = 0.61$।
\textbf{12.} \ce{O^{2-}} FCC \omterm{def:b1:crystals-metals:lattice}{जालक} पर (कोने और फलक-केंद्र),
\ce{Li+} 8 \omterm{def:b1:crystals-metals:site}{चतुष्फलकीय स्थलों} में।
\textbf{13.} $a\sqrt3/4 = \qty{200.0}{pm}$; $59 + 142 = \qty{201}{pm}$।
\textbf{14.} $\rho = 4 \times 29.88/(6.022 \times 10^{23} \times (4.619
\times 10^{-8})^3) = \qty{2.01}{g/cm^3}$, जो मापा गया मान ही है।
\textbf{15.} ज़िंक ब्लेंड: हीरा वह ज़िंक ब्लेंड है जिसमें दोनों प्रकार की स्थितियों पर
कार्बन है।
\textbf{16.} $a\sqrt3/4 = \qty{154.5}{pm}$; $\rho = 8 \times
12.01/(6.022 \times 10^{23} \times (3.567 \times 10^{-8})^3) =
\qty{3.52}{g/cm^3}$।
\textbf{17.} हीरे में केवल आधे \omterm{def:b1:crystals-metals:site}{चतुष्फलकीय स्थल} भरे हैं, और वह भी
\omterm{def:b1:crystals-metals:lattice}{जालक} के परमाणुओं जितने बड़े परमाणुओं से, जो केवल आबंधों के अनुदिश छूते हैं:
$C = 0.34$। फ़्लुओराइट में सभी स्थल भरे हैं और छोटे धनायन
बड़े ऋणायनों के बीच की ख़ाली जगह भरते हैं: $C = 0.61$।
\textbf{18.} $40.08 + 2 \times 19.00 = \qty{78.08}{g/mol}$।
\textbf{19.} $4 \times 78.08/(6.022 \times 10^{23}) = \qty{5.186e-22}{g}$।
\textbf{20.} $(5.463 \times 10^{-8})^3 = \qty{1.630e-22}{cm^3}$।
\textbf{21.} हर सौ कोष्ठिकाओं में एक \ce{F-} की कमी द्रव्यमान का $19.00/(100
\times 312.3)$, यानी 0.06\,\%, घटा देती है: तीन सार्थक अंकों पर अदृश्य (और एक
\ce{Ca^{2+}} को भी जाना होगा, अन्यथा \omterm{def:b1:crystals-metals:crystal}{क्रिस्टल} उदासीन नहीं रहेगा)।
\textbf{22.} $\rho = 5.186 \times 10^{-22}/1.630 \times 10^{-22} =
\textbf{\qty{3.18}{g/cm^3}}$, ठीक मापा गया घनत्व: $a$ के एक X-किरण मापन से
बना कोष्ठिका का मॉडल खनिज के द्रव्यमान का पूरा हिसाब देता है।
\end{solution}
